\scanpage{54}
\diagram{\begin{tikzcd}[ampersand replacement=\&]
\&X_s\arrow[r]\&Y_s\\
k(s)=k\&T_s\&
\end{tikzcd}}
\diagram{\begin{tikzcd}[ampersand replacement=\&]
\&X'_s\arrow[r,dash]\&Y'_s\\
k'\&T'_s\arrow[u]\arrow[ur]\&
\end{tikzcd}}
\margin{Nb de comp. connexes absolues de $T_s$: $\unsure{\nu|v}_s$}
\[
\left\langle\begin{aligned}
\nu&=\sup\nu_s\\
&\hspace{-1em}\left\{\begin{aligned}n&=\deg T/S\\d&=\dim.\ \text{fibres de }g\end{aligned}\right.\\
\lambda&=\deg X/Y
\end{aligned}\right.
\]
\diagram{\begin{tikzcd}[ampersand replacement=\&]
\&\deleted{\unsure{X'_s|Y'_s}}\\
\deleted{\unsure{k'}}\arrow[r,dash]\&\deleted{\unsure{T'_s}}\arrow[u]
\end{tikzcd}}
\uline{Corollaire 1 bis} Sous les conditions du corollaire 1,
\deleted{Supposons $T/S$ fini,} \deleted{Alors} et supposons que
\deleted{\unsure{$f$ soit}} \inserted{le morphisme étale} $g:X\to Y$ soit propre de degré \deleted{$\unsure{n|d}$} $\lambda$.
Alors $f_\ast(g)_s:f_\ast(X/T)_s\unsure{\mathbin{-}|\to}f_\ast(Y/T)_s$ est
propre et de degré $\lambda^{\deleted{\unsure{n}}\nu_s}$ ($\nu_s=\deg_{\mathrm{sep}}T_s/k(s)$).

On reprend le raisonnement du cor. 1,
[bloquer cor. 3 et cor. 1 en un
seul énoncé] Il suffit de prouver
qu'il y a \uline{\unsure{exactmt|exact}} $\lambda^{\nu_s}$ points dans
l'image inverse d'un pt de $f_\ast(Y/T)$.
% margin note written sideways
\margin{Sous ces conditions, un morphisme fini étale de degré $\lambda$ est transformé \unsure{en un|en un tel|en} de degré $\lambda^\nu$}

\uline{Corollaire 3} (\deleted{$T/S$ fini} \inserted{Sous les cond. \deleted{\unsure{rappellées sur la première}} du cor 1}).
Conditions équivalentes:
\diagram{\begin{tikzcd}[ampersand replacement=\&,row sep=small,column sep=small]
\text{(i)}\arrow[d]\& \nu(s)\ \text{est \inserted{loc.} constant}\\
\text{(ii)}\arrow[d]\& \begin{gathered} X\mapsto f_\ast(X/T)\ \text{transforme morphismes \deleted{finis} étales}\\ \text{\uline{finis} en morphismes étales \uline{finis}}\end{gathered}\\
\text{(iii)}\arrow[d]\& \begin{gathered}\text{Si }T'\text{ est étale fini sur }T,\ \text{alors }f_\ast(T'/T)\\ \text{est étale fini sur }S\end{gathered}\\
\text{(iv)}\arrow[uuu,bend left=24]\& f_\ast(T\amalg T)\inserted{=\underline{\operatorname{Hom}}_S(T,S\amalg S)}\quad
\text{est étale fini sur }S\unsure{\rho|\circ}
\end{tikzcd}}
% archive sheet number: 53
